\documentclass[11pt]{article}

\usepackage[margin=1in]{geometry}
\usepackage{amsmath,amssymb,amsthm,mathtools,bm}
\usepackage{enumitem}
\usepackage{booktabs}
\usepackage{microtype}
\usepackage[colorlinks=true,linkcolor=blue,citecolor=blue,urlcolor=blue]{hyperref}

\newtheorem{theorem}{Theorem}
\newtheorem{corollary}[theorem]{Corollary}
\newtheorem{proposition}[theorem]{Proposition}
\newtheorem{lemma}[theorem]{Lemma}
\newtheorem{assumption}[theorem]{Assumption}
\newtheorem{remark}[theorem]{Remark}
\newtheorem{definition}[theorem]{Definition}

\newcommand{\E}{\mathbb{E}}
\newcommand{\Pp}{\mathbb{P}}
\newcommand{\F}{\mathcal{F}}
\newcommand{\G}{\mathcal{G}}
\newcommand{\Q}{\mathcal{Q}}
\newcommand{\I}{\mathcal{I}}
\newcommand{\cE}{\mathcal{E}}
\newcommand{\1}{\mathbf{1}}
\newcommand{\R}{\mathbb{R}}
\newcommand{\eps}{\varepsilon}
\newcommand{\whp}{\widehat p}
\newcommand{\bp}{\bar p}
\newcommand{\lp}{\underline p}
\newcommand{\up}{\overline p}

\title{Delayed Bivariate Martingale DKW-Type Bounds\\
for Online Pairwise Event Probabilities}
\author{Working technical note}
\date{\today}

\begin{document}
\maketitle

\begin{abstract}
We study online estimation of a bivariate lower-orthant probability when each
observation is revealed only after a fixed delay and the observations may be
adaptively dependent.  The target is the running average of origin-conditional
bivariate CDFs,
\[
    p_k(a,b)
    =
    \Pp(B_{k,1}\le a,\,B_{k,2}\le b\mid \F_k),
\]
rather than an i.i.d. population CDF.  A delay-$d$ staggering construction
turns the delayed observations into $d$ martingale-difference streams.  On a
finite bivariate query grid, conditional Hoeffding concentration then yields a
single high-probability event that is simultaneous over grid points, residue
classes, edges, and online sample sizes.  Aggregating the staggered streams
gives the explicit delay-aware DKW-type width
\[
    \sqrt{
        \frac{d}{2N_t}
        \log\!\left(
            \frac{\pi^2 dQ N_t^2}{3\delta}
        \right)
    },
\]
where $N_t$ is the number of fully resolved origins by physical time $t$.
When $N_t\asymp t$, this is
\[
    O\!\left(
        \sqrt{
            \frac{d}{t}
            \log\!\left(\frac{dQt^2}{\delta}\right)
        }
    \right).
\]
If the conditional bivariate CDF is stable, this is a confidence band for the
common CDF.  If the CDF drifts, the same statistical term combines with an
explicit tracking-bias term; a finite-horizon moving-window version is also
derived.  We additionally study adaptive convex aggregation of several memory
scales.  A simultaneous confidence event remains valid after data-dependent
aggregation, while a residue-wise fixed-share update yields interval-wise
adaptive regret for Brier loss and therefore an oracle inequality for average
squared CDF estimation error.  A cumulative lower-orthant variation budget then
turns this oracle inequality into an explicit nonstationary estimation bound,
with the meta-learner adapting to the unknown variation scale through its choice
of memory.  The note is deliberately finite-grid: under
arbitrary adapted data, a continuum-uniform martingale DKW extension is not
automatic, and sequential rather than classical empirical-process complexity
is the relevant object.
\end{abstract}

\section{Motivation and scope}

Many online problems require estimating a probability of simultaneous events.
A useful representation is
\begin{equation}
    p_k(a,b)
    :=
    \Pp\!\left(
        B_{k,1}\le a,\;B_{k,2}\le b
        \mid \F_k
    \right),
    \label{eq:conditional-cdf}
\end{equation}
where $B_k=(B_{k,1},B_{k,2})$ is generated at origin $k$ but is only fully
observed after a delay.  In the trajectory-risk application motivating this
note, $B_k$ is a pair of critical conformal levels and
$p_k(a,b)$ is the joint failure probability of two forecast horizons at
marginal risk levels $(a,b)$.

The classical Dvoretzky--Kiefer--Wolfowitz (DKW) inequality concerns an i.i.d.
empirical CDF.  Here we instead want three features simultaneously:
\begin{enumerate}[leftmargin=*]
    \item \emph{adaptive dependence}: no independence across origins is assumed;
    \item \emph{bivariate queries}: the object of interest is a lower-orthant CDF;
    \item \emph{delayed observation}: $B_k$ is unavailable at the time its
    conditional law in \eqref{eq:conditional-cdf} is defined.
\end{enumerate}

The main result below is a finite-grid martingale analogue of a bivariate DKW
bound.  It does not require i.i.d. data, stationarity, mixing, or independence
between edges.  The price is that the theorem controls a finite collection of
lower-orthant queries rather than the entire continuum.

\section{Relation to martingale empirical-process results}

Rakhlin, Sridharan, and Tewari study the dependent empirical process
\begin{equation}
    \frac1n\sum_{i=1}^n
    \left(
        f(Z_i)-\E[f(Z_i)\mid \mathcal A_{i-1}]
    \right)
    \label{eq:rst-process}
\end{equation}
uniformly over a function class $f\in\mathcal H$, using sequential
symmetrization, sequential covering numbers, and sequential fat-shattering
complexity \cite{rakhlin2015}.  Their Lemmas 10--11 give high-probability
bounds of precisely this martingale-empirical-process form.

Yin, Agrawal, and Zeevi use these results to state a martingale version of the
univariate DKW inequality for adaptive reported values \cite{yin2022}.  Their
Lemma 1 compares the empirical CDF to the running average conditional CDF,
which is the same type of target used in this note.

There is, however, an important verification point.  Rakhlin et al. emphasize
that the relevant complexity under arbitrary adapted dependence is
\emph{sequential} complexity, and they explicitly note that classical
threshold classes need not satisfy unrestricted sequential uniform
convergence.  Consequently, extending a continuum-uniform univariate
martingale DKW statement to bivariate lower orthants cannot be justified merely
from ordinary VC dimension or ordinary fat-shattering dimension.

For this reason, the main theorem here avoids a continuum-complexity claim.  We
work on a finite query grid and prove the concentration result directly from
conditional Hoeffding inequalities.  The result is therefore self-contained
once the delayed martingale filtration is established.

Mineiro and Howard provide another complementary viewpoint: time-uniform and
value-uniform confidence bands for the \emph{running averaged conditional CDF}
under nonstationarity \cite{mineiro2023}.  This reinforces an important
interpretation used below: absent additional temporal regularity, the natural
martingale target is an average of conditional laws, not an arbitrary current
conditional law.

\section{Delayed bivariate observations}

Let $(\F_t)_{t\ge0}$ be a filtration representing information available at
physical time $t$.  Fix one pairwise process with deterministic delay
$d\in\mathbb N$.  For every origin $k\ge1$, let
\[
    B_k=(B_{k,1},B_{k,2})\in\R^2
\]
be measurable with respect to $\F_{k+d}$.  Thus the pair generated at origin
$k$ is fully observed by physical time $k+d$.

Let
\[
    \Q=\G_1\times\G_2
\]
be a finite set of lower-orthant queries, with
$Q:=|\Q|$.  In the common square-grid case,
$\G_1=\G_2=\G$ and $Q=M^2$ where $M=|\G|$.

For $g=(a,b)\in\Q$, define
\begin{align}
    X_k(g)
    &:=
    \1\{B_{k,1}\le a,\;B_{k,2}\le b\},
    \label{eq:Xkg}
    \\
    p_k(g)
    &:=
    \E[X_k(g)\mid\F_k].
    \label{eq:pkg}
\end{align}
The function $g\mapsto p_k(g)$ is the origin-$k$ conditional bivariate CDF
restricted to the query set $\Q$.

\subsection{Why delay breaks the naive martingale indexing}

The centered variable
\[
    X_k(g)-p_k(g)
\]
has mean zero conditional on $\F_k$, but it is not observed until time $k+d$.
If all origins are indexed consecutively, then the filtration immediately
preceding the reveal of $X_k(g)$ contains information strictly beyond $\F_k$.
Thus the sequence is not automatically a martingale-difference sequence with
respect to the reveal-time filtration while retaining $p_k(g)$ as the target.

The following staggering construction fixes this exactly.

\section{Staggered delayed martingale streams}

For residue $c\in\{1,\ldots,d\}$, define the origins
\begin{equation}
    k_{c,i}
    :=
    c+(i-1)d,
    \qquad i=1,2,\ldots.
    \label{eq:origin-residue}
\end{equation}
Define the staggered filtration
\begin{equation}
    \mathcal H_{c,i}
    :=
    \F_{c+id}.
    \label{eq:staggered-filtration}
\end{equation}
Then
\[
    \mathcal H_{c,i-1}
    =
    \F_{k_{c,i}},
\]
and $B_{k_{c,i}}$ is measurable with respect to
$\mathcal H_{c,i}$ because
$k_{c,i}+d=c+id$.

For $g\in\Q$, write
\begin{align}
    X_{c,i}(g)
    &:= X_{k_{c,i}}(g),
    \\
    p_{c,i}(g)
    &:= p_{k_{c,i}}(g),
    \\
    D_{c,i}(g)
    &:= X_{c,i}(g)-p_{c,i}(g).
\end{align}

\begin{lemma}[Delayed martingale difference]
\label{lem:delayed-mds}
For every residue $c$ and query $g\in\Q$,
$(D_{c,i}(g))_{i\ge1}$ is a martingale-difference sequence with respect to
$(\mathcal H_{c,i})_{i\ge0}$, and $D_{c,i}(g)\in[-1,1]$.
\end{lemma}

\begin{proof}
The variable $X_{c,i}(g)$ is $\mathcal H_{c,i}$-measurable.  Moreover,
\[
    \E[X_{c,i}(g)\mid\mathcal H_{c,i-1}]
    =
    \E[X_{k_{c,i}}(g)\mid\F_{k_{c,i}}]
    =
    p_{c,i}(g).
\]
Hence
$\E[D_{c,i}(g)\mid\mathcal H_{c,i-1}]=0$.
Since $X_{c,i}(g)\in\{0,1\}$ and $p_{c,i}(g)\in[0,1]$, the increment lies in
$[-1,1]$.
\end{proof}

The staggering is therefore a filtration device.  It is not an independence
assumption and it does not assert that different residue classes are
independent.

\section{A time-uniform finite-grid martingale DKW-type bound}

For one residue class, define the empirical and average conditional CDFs
\begin{align}
    \widehat p_{c,n}(g)
    &:=
    \frac1n\sum_{i=1}^n X_{c,i}(g),
    \label{eq:residue-empirical}
    \\
    \bar p_{c,n}(g)
    &:=
    \frac1n\sum_{i=1}^n p_{c,i}(g).
    \label{eq:residue-target}
\end{align}

\begin{theorem}[Anytime delayed bivariate grid bound]
\label{thm:single-edge-anytime}
Fix $\delta\in(0,1)$.  For $n\ge1$, define
\begin{equation}
    r_n(\delta)
    :=
    \sqrt{
        \frac{1}{2n}
        \log\!\left(
            \frac{\pi^2 dQ n^2}{3\delta}
        \right)
    }.
    \label{eq:rn-single}
\end{equation}
Then, with probability at least $1-\delta$,
\begin{equation}
    \left|
        \widehat p_{c,n}(g)-\bar p_{c,n}(g)
    \right|
    \le
    r_n(\delta)
    \label{eq:residue-band}
\end{equation}
simultaneously for every residue $c\in\{1,\ldots,d\}$, every query
$g\in\Q$, and every sample size $n\ge1$.
\end{theorem}

\begin{proof}
Fix $c$, $g$, and $n$.  Conditional Hoeffding concentration for the bounded
martingale-difference sequence in Lemma~\ref{lem:delayed-mds} gives
\[
    \Pp\!\left(
        \left|
            \widehat p_{c,n}(g)-\bar p_{c,n}(g)
        \right|>x
    \right)
    \le
    2e^{-2nx^2}.
\]
With $x=r_n(\delta)$, the right-hand side equals
\[
    \frac{6\delta}{\pi^2 dQ n^2}.
\]
Union bounding over $d$ residues and $Q$ grid queries, and then over all
$n\ge1$, gives total failure probability at most
\[
    dQ\sum_{n=1}^\infty
    \frac{6\delta}{\pi^2 dQ n^2}
    =
    \delta.
\]
\end{proof}

\begin{remark}[One confidence parameter]
The proof uses multiplicity mathematically, but the algorithm requires only one
user-specified level $\delta$.  The factors $d$ and $Q$ appear logarithmically
inside a single global radius.  There is no need to maintain separate
confidence levels $\delta_{c,g}$.
\end{remark}

\section{Aggregating residues into one bivariate CDF estimate}

At physical time $t$, let $n_{t,c}$ be the number of fully resolved origins in
residue $c$, and define
\[
    N_t:=\sum_{c=1}^d n_{t,c}.
\]
With one origin issued at every physical time and no missing observations,
$N_t=(t-d)_+$ up to the initial indexing convention.

Define the aggregate empirical CDF and its aggregate conditional target by
\begin{align}
    \widehat p_t(g)
    &:=
    \frac1{N_t}
    \sum_{c=1}^d
    \sum_{i=1}^{n_{t,c}}
    X_{c,i}(g),
    \label{eq:aggregate-empirical}
    \\
    \bar p_t(g)
    &:=
    \frac1{N_t}
    \sum_{c=1}^d
    \sum_{i=1}^{n_{t,c}}
    p_{c,i}(g).
    \label{eq:aggregate-target}
\end{align}
When $N_t=0$, the estimate is left undefined.

\begin{corollary}[Edge-level delayed CDF band]
\label{cor:aggregate-band}
On the event of Theorem~\ref{thm:single-edge-anytime}, for every $t$ with
$N_t>0$ and every $g\in\Q$,
\begin{equation}
    \left|
        \widehat p_t(g)-\bar p_t(g)
    \right|
    \le
    E_t,
    \qquad
    E_t
    :=
    \sum_{c:n_{t,c}>0}
    \frac{n_{t,c}}{N_t}
    r_{n_{t,c}}(\delta).
    \label{eq:exact-Et}
\end{equation}
Moreover,
\begin{equation}
    \boxed{
    E_t
    \le
    \sqrt{
        \frac{d}{2N_t}
        \log\!\left(
            \frac{\pi^2 dQ N_t^2}{3\delta}
        \right)
    }.
    }
    \label{eq:simple-Et}
\end{equation}
When one origin is issued at every physical time and no observations are
missing, $N_t=(t-d)_+$ up to the indexing convention.  Therefore the preceding
bound can be written explicitly as
\[
    E_t
    \le
    \sqrt{
        \frac{d}{2(t-d)}
        \left[
            2\log(t-d)
            +
            \log\!\left(\frac{\pi^2 dQ}{3\delta}\right)
        \right]
    },
    \qquad t>d,
\]
and hence, when $t\gg d$,
\[
    E_t
    =
    O\!\left(
        \sqrt{
            \frac{d}{t}
            \log\!\left(\frac{dQt^2}{\delta}\right)
        }
    \right).
\]
\end{corollary}

\begin{proof}
The first inequality follows by writing the aggregate difference as the
weighted sum of the residue-level differences.  For the second, set
\[
    L_t
    :=
    \log\!\left(
        \frac{\pi^2 dQ N_t^2}{3\delta}
    \right).
\]
Since $n_{t,c}\le N_t$,
$r_{n_{t,c}}(\delta)\le\sqrt{L_t/(2n_{t,c})}$.  Therefore
\begin{align*}
    E_t
    &\le
    \frac{\sqrt{L_t/2}}{N_t}
    \sum_{c:n_{t,c}>0}\sqrt{n_{t,c}}
    \\
    &\le
    \frac{\sqrt{L_t/2}}{N_t}
    \sqrt{d\sum_c n_{t,c}}
    =
    \sqrt{\frac{dL_t}{2N_t}},
\end{align*}
where the second inequality is Cauchy--Schwarz.
\end{proof}

\section{What distribution is being estimated?}

The martingale theorem controls
\[
    \bar p_t(g)
    =
    \frac1{N_t}\sum_{k\in\I_t} p_k(g),
\]
the running average of the origin-conditional bivariate CDFs over resolved
origins $\I_t$.  This distinction is fundamental.

\subsection{Stable conditional CDF}

\begin{assumption}[Stable conditional bivariate CDF]
\label{ass:stable-conditional}
There exists a function $p:\Q\to[0,1]$ such that
\[
    p_k(g)=p(g)
\]
for every origin $k$ and query $g\in\Q$.
\end{assumption}

This is not an i.i.d. assumption.  The realized pairs $B_k$ may be serially
dependent and may overlap in physical time.  The assumption says only that the
origin-conditional lower-orthant probabilities are invariant.

\begin{corollary}[Delayed bivariate confidence band under conditional stability]
\label{cor:stable-band}
Under Assumption~\ref{ass:stable-conditional}, with probability at least
$1-\delta$,
\begin{equation}
    \sup_{g\in\Q}
    \left|
        \widehat p_t(g)-p(g)
    \right|
    \le E_t
    \label{eq:stable-cdf-band}
\end{equation}
simultaneously for all physical times $t$ with $N_t>0$.
Consequently,
\begin{align}
    \underline p_t(g)
    &:=
    \left[\widehat p_t(g)-E_t\right]_+,
    \\
    \overline p_t(g)
    &:=
    \left[\widehat p_t(g)+E_t\right]_1
\end{align}
form simultaneous lower and upper confidence surfaces on $\Q$.
\end{corollary}

Because the same scalar $E_t$ is subtracted from every grid point, the lower
surface preserves coordinatewise monotonicity of the empirical bivariate CDF.

\subsection{Changing conditional CDFs}

If $p_k$ varies with $k$, the statistical theorem remains valid for the running
average $\bar p_t$.  To compare the empirical CDF with the current conditional
CDF $p_t$, define the deterministic tracking discrepancy
\begin{equation}
    B_t
    :=
    \sup_{g\in\Q}
    \left|p_t(g)-\bar p_t(g)\right|.
    \label{eq:tracking-Bt}
\end{equation}
Then on the same event,
\begin{equation}
    \sup_{g\in\Q}
    \left|
        \widehat p_t(g)-p_t(g)
    \right|
    \le
    E_t+B_t.
    \label{eq:current-decomp}
\end{equation}
Thus statistical martingale error and temporal tracking error separate
additively.

A convenient sufficient condition is bounded one-step variation.  Let
\begin{equation}
    v_u
    :=
    \sup_{g\in\Q}
    |p_{u+1}(g)-p_u(g)|.
    \label{eq:variation-vu}
\end{equation}
Then
\begin{equation}
    B_t
    \le
    \frac1{N_t}
    \sum_{k\in\I_t}
    \sum_{u=k}^{t-1}v_u.
    \label{eq:variation-full-history}
\end{equation}
This bound also shows why a full-history estimator may track a moving CDF
poorly: old observations accumulate temporal bias.

\section{Finite-horizon moving-window bound}

For a drifting CDF, a recent-window estimator is more natural.  This section
gives a finite-horizon version with a sharp dependence on the window size.
Fix a terminal physical horizon $T_{\max}$.  Consider any consecutive interval
of observations within a residue stream, of length $m$.

\begin{lemma}[Uniform concentration over residue intervals]
\label{lem:interval}
Let
\[
    C_0:=dQ.
\]
With probability at least $1-\delta$, every residue $c$, every query $g\in\Q$,
and every consecutive interval of residue indices contained before
$T_{\max}$ satisfy
\begin{equation}
    \left|
        \frac1m\sum_{i=a+1}^{a+m}D_{c,i}(g)
    \right|
    \le
    r_m^{\mathrm{win}},
    \qquad
    r_m^{\mathrm{win}}
    :=
    \sqrt{
        \frac1{2m}
        \log\!\left(
            \frac{C_0T_{\max}^2}{\delta}
        \right)
    }.
    \label{eq:window-radius}
\end{equation}
\end{lemma}

\begin{proof}
For any fixed interval, conditional Hoeffding gives probability at most
$2e^{-2m(r_m^{\mathrm{win}})^2}=2\delta/(C_0T_{\max}^2)$.
There are at most $T_{\max}^2/2$ index intervals in each residue stream, and
$C_0=dQ$ residue-query pairs.  A union bound gives total failure probability at
most $\delta$.
\end{proof}

Now let $\widehat p_t^{(W)}(g)$ use the $W$ most recent fully resolved origins
for the edge.  Combining the residue streams as before yields the statistical
bound
\begin{equation}
    \sup_{g\in\Q}
    \left|
        \widehat p_t^{(W)}(g)
        -
        \bar p_t^{(W)}(g)
    \right|
    \le
    \sqrt{
        \frac{d}{2W}
        \log\!\left(
            \frac{dQ T_{\max}^2}{\delta}
        \right)
    }.
    \label{eq:window-statistical}
\end{equation}

Suppose additionally that
\begin{equation}
    \sup_{g\in\Q}|p_{u+1}(g)-p_u(g)|\le\rho
    \qquad\text{for all }u.
    \label{eq:rho-drift}
\end{equation}
The newest fully resolved origin is at least $d$ time steps old.  The $W$
consecutive resolved origins have ages $d,d+1,\ldots,d+W-1$, so
\begin{equation}
    \sup_{g\in\Q}
    |p_t(g)-\bar p_t^{(W)}(g)|
    \le
    \rho\left(d+\frac{W-1}{2}\right).
    \label{eq:window-drift}
\end{equation}
Combining \eqref{eq:window-statistical} and \eqref{eq:window-drift} gives the
current-time bound
\begin{equation}
\boxed{
    \sup_{g\in\Q}
    \left|
        \widehat p_t^{(W)}(g)-p_t(g)
    \right|
    \le
    \sqrt{
        \frac{d}{2W}
        \log\!\left(
            \frac{dQ T_{\max}^2}{\delta}
        \right)
    }
    +
    \rho\left(d+\frac{W-1}{2}\right).
}
\label{eq:window-current-final}
\end{equation}
The first term is statistical uncertainty; the second is unavoidable tracking
bias under the assumed drift model.  Set
\[
    L_{\max}
    :=
    \log\!\left(\frac{dQ T_{\max}^2}{\delta}\right).
\]
For the continuous relaxation in $W$, ignoring only the integer and feasibility
constraints on the window length, the right-hand side is minimized at
\[
    W_{\star}
    =
    \left(
        \frac{dL_{\max}}{2\rho^2}
    \right)^{1/3}.
\]
At this value, the optimized bound is
\[
    \rho\left(d-\frac12\right)
    +
    \frac32
    \left(
        \frac{\rho dL_{\max}}{2}
    \right)^{1/3}.
\]
Thus a simpler order statement with the logarithmic factor shown explicitly is
\[
    O\!\left(
        \rho d
        +
        \left[
            \rho d
            \log\!\left(\frac{dQ T_{\max}^2}{\delta}\right)
        \right]^{1/3}
    \right).
\]
The additive $\rho d$ term is the irreducible staleness cost created by delayed
revelation.


\section{Adaptive aggregation over memory scales}
\label{sec:adaptive-aggregation}

The fixed-window bound above leaves a practical tuning question: which memory
length $W$ should be used when the drift rate is unknown or changes over time?
This section gives a rigorous aggregation construction.  The statistical and
online-learning guarantees are deliberately separated:
\begin{enumerate}[leftmargin=*]
    \item a single simultaneous martingale event controls \emph{every} base
    window and therefore any data-dependent convex combination of them; and
    \item a fixed-share multiplicative update chooses the mixture weights and
    enjoys interval-wise adaptive regret under delayed feedback.
\end{enumerate}
The weight update is the fixed-share update used in the MARA meta-algorithm of
Gradu, Hazan, and Minasyan \cite{gradu2023}.  We use only their multiplicative
aggregation mechanism and interval-regret argument; the control-specific
random expert-switching and action-shift machinery in that paper are not
needed here because a convex combination of CDF estimates is itself a valid
CDF estimate.

\subsection{A dictionary of window experts}

Fix a finite dictionary of window lengths
\[
    1\le W_1<\cdots<W_J,
\]
and consider origins $k$ for which all $J$ windows are available, i.e.
$k\ge d+W_J$.  Expert $j$ uses the $W_j$ most recent fully resolved origins,
\begin{equation}
    \mathcal W_{k,j}
    :=
    \{k-d-W_j+1,\ldots,k-d\},
\end{equation}
and outputs the empirical bivariate CDF
\begin{equation}
    \widehat p_{k,j}(g)
    :=
    \frac1{W_j}
    \sum_{r\in\mathcal W_{k,j}} X_r(g),
    \qquad g\in\Q.
    \label{eq:base-window-expert}
\end{equation}
Its conditional-CDF target is
\begin{equation}
    \bar p_{k,j}(g)
    :=
    \frac1{W_j}
    \sum_{r\in\mathcal W_{k,j}} p_r(g).
    \label{eq:base-window-target}
\end{equation}
Both $\widehat p_{k,j}$ and any function of previously resolved losses are
$\F_k$-measurable.  This predictability will be used below when the estimates
are scored against the delayed observation generated at origin $k$.

Let $\pi_k=(\pi_{k,1},\ldots,\pi_{k,J})$ be any random vector in the simplex
$\Delta_J$ that is $\F_k$-measurable.  Define the convex aggregate
\begin{align}
    \widehat p_k^{\mathrm{mix}}(g)
    &:=%
    \sum_{j=1}^J \pi_{k,j}\widehat p_{k,j}(g),
    \label{eq:mixture-cdf}
    \\
    \bar p_k^{\mathrm{mix}}(g)
    &:=%
    \sum_{j=1}^J \pi_{k,j}\bar p_{k,j}(g).
    \label{eq:mixture-target}
\end{align}
Because each base estimate is coordinatewise nondecreasing in $g$, the mixture
is also a bivariate CDF on the grid.

The mixture can equivalently be written as a learned weighting of observation
ages.  For $m=0,\ldots,W_J-1$, set
\begin{equation}
    \omega_{k,m}
    :=
    \sum_{j:W_j>m}
    \frac{\pi_{k,j}}{W_j}.
    \label{eq:induced-age-weights}
\end{equation}
Then $\omega_{k,m}\ge0$, $\sum_m\omega_{k,m}=1$, and
\begin{equation}
    \widehat p_k^{\mathrm{mix}}(g)
    =
    \sum_{m=0}^{W_J-1}
    \omega_{k,m}X_{k-d-m}(g).
    \label{eq:mixture-as-weighted-history}
\end{equation}
Moreover $\omega_{k,m+1}\le\omega_{k,m}$: mixing uniform windows induces a
nonincreasing recency profile automatically.

\begin{proposition}[Average age of the adaptive mixture]
\label{prop:adaptive-average-age}
The induced weights satisfy
\begin{equation}
    A_k
    :=
    \sum_{m=0}^{W_J-1}\omega_{k,m}(d+m)
    =
    \sum_{j=1}^J
    \pi_{k,j}
    \left(d+\frac{W_j-1}{2}\right).
    \label{eq:adaptive-average-age}
\end{equation}
If the one-step CDF drift obeys
\[
    \sup_{g\in\Q}|p_{u+1}(g)-p_u(g)|\le\rho
    \qquad\text{for every }u,
\]
then
\begin{equation}
    \sup_{g\in\Q}
    |\bar p_k^{\mathrm{mix}}(g)-p_k(g)|
    \le
    \rho A_k.
    \label{eq:adaptive-mixture-drift}
\end{equation}
\end{proposition}

\begin{proof}
Substituting \eqref{eq:induced-age-weights} and interchanging the finite sums,
\begin{align*}
    \sum_m\omega_{k,m}(d+m)
    &=
    \sum_{j=1}^J
    \frac{\pi_{k,j}}{W_j}
    \sum_{m=0}^{W_j-1}(d+m)
    \\
    &=
    \sum_{j=1}^J
    \pi_{k,j}
    \left(d+\frac{W_j-1}{2}\right).
\end{align*}
For every $m$ and $g$, the bounded-drift assumption and the triangle inequality
give
\[
    |p_k(g)-p_{k-d-m}(g)|
    \le \rho(d+m).
\]
Using the weighted-history representation of the target and then averaging the
last display with weights $\omega_{k,m}$ proves
\eqref{eq:adaptive-mixture-drift}.
\end{proof}

\subsection{Post-selection-safe confidence bands}

The interval-uniform event from Lemma~\ref{lem:interval} already controls every
consecutive window in every residue stream.  Consequently, selecting or mixing
window lengths after observing the past does not require another confidence
split over the $J$ experts.

Set
\begin{equation}
    L_{\max}
    :=
    \log\!\left(
        \frac{dQ T_{\max}^2}{\delta}
    \right),
    \qquad
    a_j
    :=
    \sqrt{\frac{dL_{\max}}{2W_j}}.
    \label{eq:expert-stat-radius}
\end{equation}

\begin{theorem}[Adaptive convex aggregation on one simultaneous event]
\label{thm:adaptive-mixture-band}
On the event of Lemma~\ref{lem:interval}, simultaneously for every origin $k$
for which all base windows are defined, every $g\in\Q$, and every
$\F_k$-measurable mixture vector $\pi_k\in\Delta_J$,
\begin{equation}
    \left|
        \widehat p_k^{\mathrm{mix}}(g)
        -
        \bar p_k^{\mathrm{mix}}(g)
    \right|
    \le
    A_k^{\mathrm{stat}},
    \qquad
    A_k^{\mathrm{stat}}
    :=
    \sum_{j=1}^J\pi_{k,j}a_j.
    \label{eq:adaptive-mixture-stat}
\end{equation}
Hence the statement holds with probability at least $1-\delta$ without an
additional factor $J$ inside the confidence level.

If, in addition, the bounded-drift condition with parameter $\rho$ holds, then
on the same event
\begin{equation}
\boxed{
    \sup_{g\in\Q}
    \left|
        \widehat p_k^{\mathrm{mix}}(g)-p_k(g)
    \right|
    \le
    C_k^{\mathrm{mix}}
    :=
    \sum_{j=1}^J
    \pi_{k,j}
    \left[
        \sqrt{\frac{dL_{\max}}{2W_j}}
        +
        \rho\left(d+\frac{W_j-1}{2}\right)
    \right].
}
\label{eq:adaptive-mixture-current-band}
\end{equation}
Consequently,
\begin{align}
    \underline p_k^{\mathrm{mix}}(g)
    &:=%
    \left[\widehat p_k^{\mathrm{mix}}(g)-C_k^{\mathrm{mix}}\right]_+,
    \\
    \overline p_k^{\mathrm{mix}}(g)
    &:=%
    \left[\widehat p_k^{\mathrm{mix}}(g)+C_k^{\mathrm{mix}}\right]_1
\end{align}
form simultaneous lower and upper confidence surfaces on $\Q$.
\end{theorem}

\begin{proof}
For every base window $j$, the argument leading to
\eqref{eq:window-statistical} applies to that window on the single event of
Lemma~\ref{lem:interval}; therefore
\[
    \sup_{g\in\Q}
    |\widehat p_{k,j}(g)-\bar p_{k,j}(g)|
    \le a_j
\]
simultaneously over all $k$ and $j$.  The event is already uniform over all
consecutive intervals, so introducing a finite dictionary of such windows does
not create a new union bound.

For any realized $\pi_k\in\Delta_J$,
\begin{align*}
    \left|
        \widehat p_k^{\mathrm{mix}}(g)
        -
        \bar p_k^{\mathrm{mix}}(g)
    \right|
    &\le
    \sum_{j=1}^J
    \pi_{k,j}
    |\widehat p_{k,j}(g)-\bar p_{k,j}(g)|
    \\
    &\le
    \sum_{j=1}^J\pi_{k,j}a_j.
\end{align*}
This argument is pathwise on the simultaneous event, which is why the weights
may be data dependent.  Finally,
\[
    |\widehat p_k^{\mathrm{mix}}(g)-p_k(g)|
    \le
    |\widehat p_k^{\mathrm{mix}}(g)-\bar p_k^{\mathrm{mix}}(g)|
    +
    |\bar p_k^{\mathrm{mix}}(g)-p_k(g)|,
\]
and Proposition~\ref{prop:adaptive-average-age} bounds the second term by
$\rho A_k$.  Substituting \eqref{eq:adaptive-average-age} proves
\eqref{eq:adaptive-mixture-current-band}.
\end{proof}

\begin{remark}[What adaptivity costs in the confidence radius]
The confidence event itself pays no extra $\log J$ factor because
Lemma~\ref{lem:interval} is uniform over all admissible intervals before the
weights are chosen.  The price of learning among $J$ memory scales appears
instead in the adaptive-regret term below.  The radius
$A_k^{\mathrm{stat}}=\sum_j\pi_{k,j}a_j$ can be conservative compared with a
concentration inequality tailored to one deterministic weighted average, but
it remains valid under arbitrary data-dependent post-selection of the mixture
weights.
\end{remark}

\subsection{Fixed-share learning under delayed feedback}

We next specify how to learn $\pi_k$.  For each origin $k$ and expert $j$, store
the prediction \eqref{eq:base-window-expert} made at origin $k$.  When $B_k$ is
revealed at physical time $k+d$, compute the integrated Brier loss
\begin{equation}
    \ell_{k,j}
    :=
    \frac1Q
    \sum_{g\in\Q}
    \left(
        X_k(g)-\widehat p_{k,j}(g)
    \right)^2
    \in[0,1].
    \label{eq:expert-brier-loss}
\end{equation}
For a mixture prediction, define analogously
\begin{equation}
    \ell_k^{\mathrm{mix}}
    :=
    \frac1Q
    \sum_{g\in\Q}
    \left(
        X_k(g)-\widehat p_k^{\mathrm{mix}}(g)
    \right)^2.
    \label{eq:mixture-brier-loss}
\end{equation}

A single global expert learner would receive these losses with delay $d$.
Instead, use the same residue decomposition that produced the martingale
filtration.  For residue $c\in\{1,\ldots,d\}$, write its eligible origins as
\[
    k_{c,1}<k_{c,2}<\cdots,
    \qquad
    k_{c,i+1}=k_{c,i}+d.
\]
The outcome from $k_{c,i}$ is available at $k_{c,i}+d=k_{c,i+1}$, before the
next prediction in that residue stream.  Thus each residue stream is a
standard full-information expert problem in its own meta-time index.

Let
\[
    n_{\max}:=\left\lceil\frac{T_{\max}}{d}\right\rceil,
    \qquad
    \sigma:=\frac1{n_{\max}},
    \qquad
    0<\eta\le\frac12,
\]
and assume $n_{\max}\ge2$.  For every residue, initialize $w_{c,1,j}=1$.
At meta-round $i$, use
\begin{equation}
    \pi_{k_{c,i},j}
    =
    \frac{w_{c,i,j}}{\sum_{\ell=1}^J w_{c,i,\ell}},
    \label{eq:fixed-share-probabilities}
\end{equation}
form the deterministic mixture \eqref{eq:mixture-cdf}, and, when the loss
vector from that origin is revealed, update
\begin{align}
    \widetilde w_{c,i+1,j}
    &:=%
    w_{c,i,j}
    \exp(-\eta\ell_{k_{c,i},j}),
    \label{eq:fixed-share-multiplicative}
    \\
    w_{c,i+1,j}
    &:=%
    (1-\sigma)\widetilde w_{c,i+1,j}
    +
    \frac{\sigma}{J}
    \sum_{\ell=1}^J
    \widetilde w_{c,i+1,\ell}.
    \label{eq:fixed-share-share}
\end{align}
This is the fixed-share multiplicative update used in Algorithm~2 of
Gradu, Hazan, and Minasyan \cite{gradu2023}.  Their random expert-selection
step is unnecessary here because we can output the convex mixture directly.

\begin{theorem}[Residue-wise adaptive regret for the CDF mixture]
\label{thm:residue-adaptive-regret}
Fix a residue $c$ and any contiguous meta-time interval
$I_c=[r,s]$.  Under updates
\eqref{eq:fixed-share-multiplicative}--\eqref{eq:fixed-share-share},
\begin{equation}
\boxed{
    (1-\eta)
    \sum_{i=r}^s
    \ell_{k_{c,i}}^{\mathrm{mix}}
    \le
    \min_{j\in[J]}
    \sum_{i=r}^s
    \ell_{k_{c,i},j}
    +
    \frac{\log(n_{\max}J)+2}{\eta}.
}
\label{eq:residue-adaptive-regret}
\end{equation}
Thus the learned mixture competes, on every local interval of a residue
stream, with the best fixed memory scale on that interval.
\end{theorem}

\begin{proof}
This is a direct specialization of the fixed-share potential argument in
Lemma~A.2 of Gradu, Hazan, and Minasyan \cite{gradu2023}; we reproduce the
steps to make the constants and the mapping to the CDF problem explicit.

Suppress the residue index and write
$W_i:=\sum_j w_{i,j}$ and
$\pi_{i,j}:=w_{i,j}/W_i$.  The sharing step preserves total weight, so
\begin{equation}
    \frac{W_{i+1}}{W_i}
    =
    \sum_{j=1}^J
    \pi_{i,j}e^{-\eta\ell_{i,j}}.
    \label{eq:weight-potential-ratio}
\end{equation}
Because $0\le\ell_{i,j}\le1$ and $0<\eta\le1/2$,
$e^{-x}\le 1-x+x^2$ for $x\in[0,1]$, and
$\ell_{i,j}^2\le\ell_{i,j}$.  Hence, with
$m_i:=\sum_j\pi_{i,j}\ell_{i,j}$,
\begin{align*}
    \frac{W_{i+1}}{W_i}
    &\le
    1-\eta m_i+\eta^2m_i
    \\
    &\le
    \exp\{-\eta(1-\eta)m_i\}.
\end{align*}
Multiplying over $i=r,\ldots,s$ gives
\begin{equation}
    \log\frac{W_{s+1}}{W_r}
    \le
    -\eta(1-\eta)
    \sum_{i=r}^s m_i.
    \label{eq:potential-upper}
\end{equation}

For any fixed expert $j$, the sharing update implies
\[
    w_{i+1,j}
    \ge
    (1-\sigma)w_{i,j}e^{-\eta\ell_{i,j}}.
\]
Iterating over $I_c$ and using $W_{s+1}\ge w_{s+1,j}$ yields
\begin{equation}
    \log\frac{W_{s+1}}{W_r}
    \ge
    \log\pi_{r,j}
    +|I_c|\log(1-\sigma)
    -\eta\sum_{i=r}^s\ell_{i,j}.
    \label{eq:potential-lower}
\end{equation}
After every share step, $\pi_{i,j}\ge\sigma/J$, while the initialization is
uniform, so $\pi_{r,j}\ge\sigma/J$.  With
$\sigma=1/n_{\max}$, $|I_c|\le n_{\max}$, and
$\log(1-\sigma)\ge-2\sigma$ for $\sigma\le1/2$, the right-hand side of
\eqref{eq:potential-lower} is at least
\[
    -\log(n_{\max}J)-2
    -\eta\sum_{i=r}^s\ell_{i,j}.
\]
Combining with \eqref{eq:potential-upper}, dividing by $\eta$, and minimizing
over $j$ gives
\begin{equation}
    (1-\eta)
    \sum_{i=r}^s m_i
    \le
    \min_j\sum_{i=r}^s\ell_{i,j}
    +
    \frac{\log(n_{\max}J)+2}{\eta}.
    \label{eq:weighted-expert-adreg}
\end{equation}
Finally, squared loss is convex in the prediction.  Therefore Jensen's
inequality gives, at every origin,
\[
    \ell_{k_{c,i}}^{\mathrm{mix}}
    \le
    \sum_j\pi_{i,j}\ell_{i,j}
    =m_i.
\]
Substitution into \eqref{eq:weighted-expert-adreg} proves the theorem.
\end{proof}

\begin{corollary}[Adaptive regret on a physical-time interval]
\label{cor:physical-adaptive-regret}
Let $I=[r,s]$ be any physical-time interval after the common burn-in and let
$d_I$ denote the number of residue classes represented in $I$; hence
$d_I\le\min\{d,|I|\}$.  Then
\begin{equation}
\boxed{
    (1-\eta)
    \sum_{k\in I}
    \ell_k^{\mathrm{mix}}
    \le
    \min_{j\in[J]}
    \sum_{k\in I}
    \ell_{k,j}
    +
    \frac{d_I\,[\log(n_{\max}J)+2]}{\eta}.
}
\label{eq:physical-adaptive-regret}
\end{equation}
\end{corollary}

\begin{proof}
The intersection of a physical interval with any residue class is a contiguous
meta-time interval.  Apply Theorem~\ref{thm:residue-adaptive-regret} to every
nonempty residue intersection and sum.  This gives the sum, over residues, of
their individual best-expert losses plus $d_I$ copies of the complexity term.
For every common expert $j$, the sum of the residue-wise minima is no larger
than that expert's loss over all of $I$; minimizing over the common $j$ gives
\eqref{eq:physical-adaptive-regret}.
\end{proof}

\subsection{From Brier adaptive regret to CDF estimation error}

The Brier loss has a direct statistical interpretation in this problem.  For
any $\F_k$-measurable CDF prediction $q:\Q\to[0,1]$, define its grid-averaged
squared estimation error
\begin{equation}
    \mathcal R_k(q)
    :=
    \frac1Q
    \sum_{g\in\Q}
    (q(g)-p_k(g))^2.
    \label{eq:grid-mse-risk}
\end{equation}
Also define the irreducible Bernoulli variance term
\begin{equation}
    \nu_k
    :=
    \frac1Q
    \sum_{g\in\Q}
    p_k(g)(1-p_k(g))
    \le\frac14.
    \label{eq:brier-noise}
\end{equation}

\begin{remark}[Interpretation as excess risk relative to the true CDF]
\label{rem:truth-oracle-risk}
The risk
\[
    \mathcal R_k(q)
    =
    \frac1Q
    \sum_{g\in\Q}
    \bigl(q(g)-p_k(g)\bigr)^2
\]
has a natural clairvoyant benchmark.  An oracle that knows the current
conditional CDF $p_k$ exactly would choose $q=p_k$ and incur zero risk:
\[
    \mathcal R_k(p_k)=0.
\]
Consequently,
\[
    \sum_{k\in I}
    \E\!\left[
        \mathcal R_k(\widehat p_k^{\mathrm{mix}})
    \right]
\]
is the cumulative expected excess CDF-estimation risk relative to a
clairvoyant oracle that knows $p_k$ at every origin.

The oracle inequalities developed below therefore admit two complementary
interpretations.  First, the adaptive mixture competes with the best fixed
memory scale in the chosen window dictionary.  Second, after bounding the
risk of that best window, the resulting bound directly controls the
aggregate's absolute gap to the zero-risk truth oracle.  In this decomposition,
the best-window term captures the intrinsic statistical and tracking
difficulty of estimating the time-varying CDF, while the remaining
meta-learning terms quantify the additional price of adaptively learning an
appropriate memory scale.
\end{remark}

\begin{lemma}[Conditional Brier decomposition]
\label{lem:brier-decomposition}
For every $\F_k$-measurable prediction $q$,
\begin{equation}
    \E\!\left[
        \frac1Q\sum_{g\in\Q}(X_k(g)-q(g))^2
        \;\middle|\;\F_k
    \right]
    =
    \nu_k+\mathcal R_k(q).
    \label{eq:brier-decomposition}
\end{equation}
\end{lemma}

\begin{proof}
For a fixed $g$, $X_k(g)\in\{0,1\}$ and
$\E[X_k(g)\mid\F_k]=p_k(g)$.  Expanding around the conditional mean,
\begin{align*}
    \E[(X_k(g)-q(g))^2\mid\F_k]
    &=
    \operatorname{Var}(X_k(g)\mid\F_k)
    +(p_k(g)-q(g))^2
    \\
    &=
    p_k(g)(1-p_k(g))+(p_k(g)-q(g))^2.
\end{align*}
Average over $g\in\Q$.
\end{proof}

\begin{theorem}[Local oracle inequality for adaptive CDF estimation]
\label{thm:adaptive-cdf-oracle}
For every physical interval $I=[r,s]$ after burn-in, the residue-wise
fixed-share aggregate satisfies
\begin{equation}
\boxed{
    \sum_{k\in I}
    \E\!\left[
        \mathcal R_k(\widehat p_k^{\mathrm{mix}})
    \right]
    \le
    \frac1{1-\eta}
    \left
    \{
        \min_{j\in[J]}
        \sum_{k\in I}
        \E\!\left[
            \mathcal R_k(\widehat p_{k,j})
        \right]
        +
        \frac{d_I\,[\log(n_{\max}J)+2]}{\eta}
        +
        \frac{\eta|I|}{4}
    \right\}.
}
\label{eq:adaptive-cdf-oracle}
\end{equation}
Thus, on every local interval, the aggregate tracks the best memory scale in
expected grid-averaged squared CDF error, up to an explicit adaptive-regret
penalty.  No drift parameter $\rho$ is required to run the weighting scheme.
\end{theorem}

\begin{proof}
Corollary~\ref{cor:physical-adaptive-regret} is a pathwise inequality in the
realized Brier losses.  Take expectations of both sides.  Since
$\E[\min_j Y_j]\le\min_j\E[Y_j]$,
\begin{align*}
    (1-\eta)
    \sum_{k\in I}\E[\ell_k^{\mathrm{mix}}]
    &\le
    \min_j
    \sum_{k\in I}\E[\ell_{k,j}]
    +
    \frac{d_I[\log(n_{\max}J)+2]}{\eta}.
\end{align*}
Every base prediction and the mixture prediction are $\F_k$-measurable by
construction.  Applying Lemma~\ref{lem:brier-decomposition} to each term gives
\begin{align*}
    (1-\eta)
    \sum_{k\in I}
    \E[\nu_k+\mathcal R_k(\widehat p_k^{\mathrm{mix}})]
    &\le
    \sum_{k\in I}\E[\nu_k]
    +
    \min_j
    \sum_{k\in I}
    \E[\mathcal R_k(\widehat p_{k,j})]
    +
    \frac{d_I[\log(n_{\max}J)+2]}{\eta}.
\end{align*}
Move the common variance contribution on the left to the right.  The remainder
is $\eta\sum_{k\in I}\E[\nu_k]\le\eta|I|/4$ by
\eqref{eq:brier-noise}.  Dividing by $1-\eta$ proves
\eqref{eq:adaptive-cdf-oracle}.
\end{proof}


\subsection{Variation-budget view of nonstationarity}
\label{subsec:variation-budget}

The pointwise drift envelope $v_u\le\rho$ is convenient for a certified
current-CDF band, but it can be unnecessarily restrictive for regret analysis:
it rules out a rare large jump even when the process is otherwise stable.  A
more flexible description is the cumulative path variation
\begin{equation}
    v_u
    :=
    \sup_{g\in\Q}|p_{u+1}(g)-p_u(g)|,
    \qquad
    V_{a:b}
    :=
    \sum_{u=a}^{b-1} v_u,
    \quad a<b,
    \label{eq:path-variation}
\end{equation}
with $V_{a:b}:=0$ when $b\le a$.  This is a lower-orthant CDF variation
budget, not total variation distance between the full laws.  The use of a
cumulative variation budget to quantify nonstationarity is standard in online
and stochastic optimization; see, for example, Besbes, Gur, and Zeevi
\cite{besbes2015}.  Their variation metric and optimization problem are
different from ours, so none of their minimax regret rates are invoked below.
Here the variation budget is used only to control the tracking error of a
delayed CDF estimator.

The bounded-drift model is the special case
\[
    V_{a:b}\le \rho(b-a).
\]
The cumulative formulation also permits abrupt changes followed by long stable
periods.

\begin{proposition}[Tracking error under an arbitrary variation path]
\label{prop:mixture-variation-path}
For the adaptive mixture in \eqref{eq:mixture-as-weighted-history}, define
\[
    B_k^{\mathrm{mix}}
    :=
    \sup_{g\in\Q}
    |\bar p_k^{\mathrm{mix}}(g)-p_k(g)|.
\]
Then, pathwise,
\begin{equation}
\boxed{
    B_k^{\mathrm{mix}}
    \le
    \sum_{m=0}^{W_J-1}
    \omega_{k,m}
    V_{k-d-m:k}.
}
\label{eq:mixture-variation-tracking}
\end{equation}
For a single window expert $j$,
\begin{equation}
    B_{k,j}
    :=
    \sup_{g\in\Q}
    |\bar p_{k,j}(g)-p_k(g)|
    \le
    \frac1{W_j}
    \sum_{m=0}^{W_j-1}
    V_{k-d-m:k}.
    \label{eq:window-variation-tracking}
\end{equation}
Consequently, if $v_u\le\rho$ for every $u$, then
\eqref{eq:mixture-variation-tracking} reduces to
$B_k^{\mathrm{mix}}\le\rho A_k$, recovering
Proposition~\ref{prop:adaptive-average-age}.
\end{proposition}

\begin{proof}
For every $m$ and $g$, repeated triangle inequality gives
\[
    |p_k(g)-p_{k-d-m}(g)|
    \le
    \sum_{u=k-d-m}^{k-1}
    \sup_{h\in\Q}|p_{u+1}(h)-p_u(h)|
    =
    V_{k-d-m:k}.
\]
The target mixture has the same age weights as
\eqref{eq:mixture-as-weighted-history}, namely
$\bar p_k^{\mathrm{mix}}(g)=\sum_m\omega_{k,m}p_{k-d-m}(g)$.
Taking the weighted average of the preceding inequality proves
\eqref{eq:mixture-variation-tracking}.  Uniform weights
$\omega_{k,m}=1/W_j$ for $m=0,\ldots,W_j-1$ give
\eqref{eq:window-variation-tracking}.  Finally,
$V_{k-d-m:k}\le\rho(d+m)$ under $v_u\le\rho$; averaging over $m$ gives
$\rho A_k$ by \eqref{eq:adaptive-average-age}.
\end{proof}

The next lemma quantifies the stochastic part of a fixed-window estimator in
mean square.  It uses only the staggered martingale-difference construction;
no mixing assumption is needed.

\begin{lemma}[Mean-square fluctuation of one window expert]
\label{lem:window-mse-fluctuation}
For every base window $j$ and eligible origin $k$,
\begin{equation}
    \E\!\left[
        \frac1Q
        \sum_{g\in\Q}
        \left(
            \widehat p_{k,j}(g)-\bar p_{k,j}(g)
        \right)^2
    \right]
    \le
    \frac{d}{4W_j}.
    \label{eq:window-mse-fluctuation}
\end{equation}
\end{lemma}

\begin{proof}
Fix $g\in\Q$ and partition the $W_j$ origins in
$\mathcal W_{k,j}$ by their residue modulo $d$.  For residue $c$, let
\[
    S_c(g)
    :=
    \sum_{r\in\mathcal W_{k,j}:\,r\equiv c\!\!\pmod d}
    \bigl(X_r(g)-p_r(g)\bigr).
\]
Then
\[
    \widehat p_{k,j}(g)-\bar p_{k,j}(g)
    =
    \frac1{W_j}\sum_{c=1}^d S_c(g).
\]
By Cauchy--Schwarz,
$\left(\sum_c S_c\right)^2\le d\sum_c S_c^2$.  Within each residue stream,
the summands are martingale differences with respect to the staggered
filtration.  Distinct increments in the same stream are therefore orthogonal
in $L^2$, and for a Bernoulli increment
\[
    \E[(X_r(g)-p_r(g))^2]
    =
    \E[p_r(g)(1-p_r(g))]
    \le \frac14.
\]
If $n_c$ window observations lie in residue $c$, then
$\E[S_c(g)^2]\le n_c/4$.  Since $\sum_c n_c=W_j$,
\[
    \E\!\left[
        (\widehat p_{k,j}(g)-\bar p_{k,j}(g))^2
    \right]
    \le
    \frac{d}{4W_j}.
\]
Averaging over $g\in\Q$ proves the lemma.
\end{proof}

\begin{lemma}[Cumulative tracking cost of a fixed window]
\label{lem:cumulative-window-variation}
Let $I=[r,s]$ be a physical-time interval after burn-in and set
$n_I:=|I|$.  For expert $j$, define the enlarged-interval variation
\begin{equation}
    \mathcal V_{I,j}
    :=
    V_{r-d-W_j+1:s}
    =
    \sum_{u=r-d-W_j+1}^{s-1}v_u,
    \label{eq:expanded-variation-budget}
\end{equation}
and its average age
\[
    A_j
    :=
    d+\frac{W_j-1}{2}.
\]
Then, pathwise,
\begin{equation}
\boxed{
    \sum_{k\in I} B_{k,j}
    \le
    A_j\,\mathcal V_{I,j}.
}
\label{eq:cumulative-window-tracking}
\end{equation}
\end{lemma}

\begin{proof}
By \eqref{eq:window-variation-tracking},
\begin{align*}
    \sum_{k=r}^s B_{k,j}
    &\le
    \frac1{W_j}
    \sum_{m=0}^{W_j-1}
    \sum_{k=r}^s
    \sum_{u=k-d-m}^{k-1}v_u.
\end{align*}
Fix $m$.  A particular increment $v_u$ can appear in the inner sum only for
$k\in\{u+1,\ldots,u+d+m\}$, hence at most $d+m$ times.  All such increments
lie inside $[r-d-W_j+1,s-1]$.  Therefore
\[
    \sum_{k=r}^s
    \sum_{u=k-d-m}^{k-1}v_u
    \le
    (d+m)\,\mathcal V_{I,j}.
\]
Averaging over $m=0,\ldots,W_j-1$ yields
\[
    \sum_{k\in I}B_{k,j}
    \le
    \left[
        \frac1{W_j}\sum_{m=0}^{W_j-1}(d+m)
    \right]
    \mathcal V_{I,j}
    =
    A_j\mathcal V_{I,j}.
\]
\end{proof}

\begin{theorem}[Variation-budget oracle bound for the adaptive estimator]
\label{thm:variation-budget-oracle}
Let $I=[r,s]$ be any physical interval after burn-in, with $n_I=|I|$.
For each expert $j$, let $\mathcal V_{I,j}$ be defined by
\eqref{eq:expanded-variation-budget}.  Then the residue-wise fixed-share
aggregate satisfies
\begin{equation}
\boxed{
\begin{aligned}
    \sum_{k\in I}
    \E\!\left[
        \mathcal R_k(\widehat p_k^{\mathrm{mix}})
    \right]
    \le
    \frac1{1-\eta}
    \Bigg[
        \min_{j\in[J]}
        \left\{
            \frac{n_I d}{2W_j}
            +
            2A_j\,\E[\mathcal V_{I,j}]
        \right\}
        +
        \frac{d_I[\log(n_{\max}J)+2]}{\eta}
        +
        \frac{\eta n_I}{4}
    \Bigg].
\end{aligned}
}
\label{eq:variation-budget-oracle}
\end{equation}
If a deterministic pathwise budget
$\mathcal V_{I,j}\le V_I$ holds for every candidate window, then
$\E[\mathcal V_{I,j}]$ may be replaced by $V_I$.
The meta-learner does not need to know $V_I$.
\end{theorem}

\begin{proof}
For any $g$,
\[
    \widehat p_{k,j}(g)-p_k(g)
    =
    \bigl(\widehat p_{k,j}(g)-\bar p_{k,j}(g)\bigr)
    +
    \bigl(\bar p_{k,j}(g)-p_k(g)\bigr).
\]
Using $(x+y)^2\le2x^2+2y^2$, averaging over $g$, taking expectations, and
applying Lemma~\ref{lem:window-mse-fluctuation} gives
\begin{equation}
    \E[\mathcal R_k(\widehat p_{k,j})]
    \le
    \frac{d}{2W_j}
    +
    2\E[B_{k,j}^2].
    \label{eq:base-risk-prevariation}
\end{equation}
Because two CDF values differ by at most one, $0\le B_{k,j}\le1$, so
$B_{k,j}^2\le B_{k,j}$.  Summing
\eqref{eq:base-risk-prevariation} over $k\in I$ and using
Lemma~\ref{lem:cumulative-window-variation} yields
\begin{equation}
    \sum_{k\in I}
    \E[\mathcal R_k(\widehat p_{k,j})]
    \le
    \frac{n_I d}{2W_j}
    +
    2A_j\E[\mathcal V_{I,j}].
    \label{eq:base-risk-variation}
\end{equation}
Substitute \eqref{eq:base-risk-variation} into the local oracle inequality
\eqref{eq:adaptive-cdf-oracle}.  The latter rests on the residue-wise
fixed-share adaptive-regret argument of Gradu, Hazan, and Minasyan
\cite{gradu2023}, specialized and proved above.  This proves
\eqref{eq:variation-budget-oracle}.  The deterministic-budget statement is
immediate.
\end{proof}

\begin{corollary}[Rate suggested by a common variation budget]
\label{cor:variation-budget-rate}
Suppose $\mathcal V_{I,j}\le V_I$ for every candidate window and that the
window dictionary contains a value within a constant factor of
\[
    W_I^*
    :=
    \sqrt{\frac{n_I d}{2V_I}},
    \label{eq:variation-optimal-window}
\]
whenever $0<V_I<\infty$ and this value lies in the admissible window range.
Then the best-window term in \eqref{eq:variation-budget-oracle} is, up to a
universal constant from discretizing the window dictionary,
\begin{equation}
    O\!\left(
        \sqrt{n_I d V_I}
        +
        dV_I
    \right).
    \label{eq:variation-rate}
\end{equation}
For a geometric dictionary $W_j=2^{j-1}$, the constant-factor coverage of
$W_I^*$ is automatic away from the endpoints.  Thus the meta-learner adapts to
the unknown variation scale through its competition with the best memory
length; no value of $V_I$ or $\rho$ is required by the update itself.
\end{corollary}

\begin{proof}
Under the common budget, the window-dependent part of
\eqref{eq:variation-budget-oracle} is
\[
    \frac{n_I d}{2W}
    +
    2\left(d+\frac{W-1}{2}\right)V_I
    =
    \frac{n_I d}{2W}
    +
    WV_I
    +(2d-1)V_I.
\]
Ignoring the additive term independent of $W$, the continuous minimizer of
$\frac{n_Id}{2W}+WV_I$ is
$W_I^*=\sqrt{n_Id/(2V_I)}$, and the two $W$-dependent terms sum to
$\sqrt{2n_IdV_I}$.  Choosing a dictionary element within a constant factor of
$W_I^*$ changes this quantity by only a universal constant.  Adding the
$(2d-1)V_I$ term proves \eqref{eq:variation-rate}.
\end{proof}

\begin{remark}[When sublinear estimation error is possible]
For fixed delay $d$, the variation-dependent contribution in
\eqref{eq:variation-rate} is sublinear in $n_I$ whenever $V_I=o(n_I)$.
This parallels the general role of sublinear variation budgets in
nonstationary regret theory, but should not be read as an application of the
minimax results of Besbes, Gur, and Zeevi \cite{besbes2015}: our target,
feedback structure, loss, and variation metric are different.  If the
variation budget is linear in the interval length, the bound permits linear
cumulative estimation error, which is the expected qualitative behavior in a
persistently changing environment.
\end{remark}

\begin{remark}[Interpretation of the two guarantees]
Theorem~\ref{thm:adaptive-mixture-band} and
Theorem~\ref{thm:adaptive-cdf-oracle} answer different questions.  The first is
a high-probability $\ell_\infty(\Q)$ confidence statement and still requires a
temporal regularity condition, such as a known $\rho$, if the target is the
current CDF $p_k$.  The second is an adaptive prediction guarantee: the
fixed-share weights need not know $\rho$ and compete locally with the best
window expert in expected squared CDF error.  Thus one may use the same
adaptive mixture either in a certified mode, when a valid drift envelope is
available, or as a plug-in estimator whose tracking quality is controlled by
adaptive regret.
\end{remark}

\begin{remark}[Choice of the memory dictionary]
A convenient implementation uses geometric windows, for example
$W_j=2^{j-1}$ up to the largest feasible window.  Then $J=O(\log T_{\max})$,
so the expert-selection complexity enters the regret bound only through
$\log J$.  Under the bounded-drift model, writing
$L_{\max}=\log(dQ T_{\max}^2/\delta)$, the continuous-relaxation optimizer of
the fixed-window confidence radius is
$W_{\star}=(dL_{\max}/(2\rho^2))^{1/3}$ as derived above.  A geometric dictionary contains a
window within a constant factor of that scale without requiring the algorithm
to know $\rho$ in advance; Theorem~\ref{thm:adaptive-cdf-oracle} explains how
the Brier-weighted aggregate can track whichever available scale predicts best
on each local interval.
\end{remark}

\section{Simultaneous control of many edges}

For applications with a set $\cE$ of pairwise processes, let edge $e$ have
delay $d_e$ and query set $\Q_e$ of size $Q_e$.  Define the global complexity
count
\begin{equation}
    C
    :=
    \sum_{e\in\cE} d_e Q_e.
    \label{eq:global-C}
\end{equation}
The proof of Theorem~\ref{thm:single-edge-anytime} gives the following result
without requiring independence across edges.

\begin{corollary}[One global confidence event over edges]
\label{cor:multi-edge}
For
\begin{equation}
    r_n^{\mathrm{glob}}(\delta)
    :=
    \sqrt{
        \frac1{2n}
        \log\!\left(
            \frac{\pi^2 C n^2}{3\delta}
        \right)
    },
    \label{eq:global-rn}
\end{equation}
with probability at least $1-\delta$, the residue-level bounds hold
simultaneously for every edge $e$, every one of its $d_e$ residues, every query
$g\in\Q_e$, and every sample size $n\ge1$.

Consequently, if $N_{t,e}$ edge-$e$ observations have resolved by time $t$,
then
\begin{equation}
\boxed{
    \sup_{g\in\Q_e}
    \left|
        \widehat p_{t,e}(g)-\bar p_{t,e}(g)
    \right|
    \le
    \sqrt{
        \frac{d_e}{2N_{t,e}}
        \log\!\left(
            \frac{\pi^2 C N_{t,e}^2}{3\delta}
        \right)
    }
}
\label{eq:multi-edge-rate}
\end{equation}
simultaneously over all edges and times.
\end{corollary}

This form uses one global confidence level $\delta$.  The multiplicity over
edges, delays, and bivariate grid points enters only through the logarithm of
$C$.

\section{Implementation}

The theorem is compatible with a simple streaming implementation.
For each edge $e$:
\begin{enumerate}[leftmargin=*]
    \item maintain an $M\times M$ histogram of resolved pairs $B_k$;
    \item increment one histogram cell whenever a new pair resolves;
    \item obtain the empirical bivariate CDF on the grid by a two-dimensional
    cumulative sum;
    \item maintain only the residue sample counts $n_{t,e,c}$, not separate
    histograms for each residue;
    \item compute the exact statistical width
    \[
        E_{t,e}
        =
        \sum_c
        \frac{n_{t,e,c}}{N_{t,e}}
        r_{n_{t,e,c}}^{\mathrm{glob}}(\delta),
    \]
    or the simpler upper bound in \eqref{eq:multi-edge-rate};
    \item return either the plug-in surface $\widehat p_{t,e}$ or the one-sided
    lower surface
    \[
        \underline p_{t,e}(g)
        =
        [\widehat p_{t,e}(g)-E_{t,e}]_+.
    \]
\end{enumerate}

Because every resolved bivariate observation determines all lower-orthant
indicators simultaneously, a single pair updates the entire empirical CDF.
The grid size affects storage and numerical resolution, while its statistical
cost is only logarithmic through $Q=M^2$.

With a dense graph of pairwise processes, storage is
$O(|\cE|M^2)$ for the histograms plus $O(\sum_e d_e)$ residue counts.  Both the
histogram updates and cumulative sums are amenable to batched tensor/GPU
implementations.

\section{Why the finite grid is theoretically useful}

In the i.i.d. setting, bivariate empirical CDFs admit continuum-uniform
Glivenko--Cantelli/DKW-type theory.  Under arbitrary adapted dependence, the
correct complexity notion changes.  Rakhlin et al. show that sequential
uniform convergence is characterized by sequential complexity measures, and
explicitly note that even the ordinary one-dimensional threshold class can
fail unrestricted sequential uniform convergence \cite{rakhlin2015}.

This creates a genuine obstacle to a completely distribution-free statement
of the form
\[
    \sup_{a,b\in\R}
    \left|
        \widehat p_t(a,b)-\bar p_t(a,b)
    \right|
    \le \epsilon_t
\]
under arbitrary adaptation.  Since bivariate lower orthants contain
one-dimensional threshold subclasses, a continuum extension needs additional
structure or a more specialized argument.

The finite grid removes this ambiguity: it creates a finite function class for
which the martingale concentration argument is immediate and fully explicit.
For online allocation applications this is often sufficient because the same
grid is already used numerically.  A continuum extension under restricted
adversaries, mixing, stability, or instance-dependent sequential complexity is
a separate theoretical direction.

\section{Discussion and open questions}

The note establishes a conservative but transparent delayed bivariate
martingale CDF theory.  Several improvements are possible.

\begin{enumerate}[leftmargin=*]
    \item \textbf{Sharper treatment of delay.}
    The $\sqrt{d/t}$ rate is produced by the staggered proof.  A specialized
    delayed-martingale argument may combine all origins more efficiently and
    reduce the delay penalty.  The present theorem should therefore be viewed
    as a rigorous baseline, not a minimax statement.

    \item \textbf{Continuum-uniform lower orthants.}
    It is worth identifying assumptions under which the bivariate lower-orthant
    class has tractable sequential complexity.  Such a result could remove the
    explicit grid from the theorem.

    \item \textbf{Time-uniform moving windows.}
    The moving-window result above is uniform over windows only up to a fixed
    terminal horizon $T_{\max}$.  Confidence-sequence stitching could replace
    this with a fully anytime windowed result.

    \item \textbf{Data-dependent widths.}
    Empirical-Bernstein or variance-adaptive martingale bounds may be much
    tighter when joint failures are rare, as is typical in low-miscoverage
    applications.

    \item \textbf{Sharper adaptive aggregation.}
    Section~\ref{sec:adaptive-aggregation} uses residue-wise fixed-share
    aggregation because it converts the physical delay into ordinary
    full-information meta-rounds.  A delayed-experts analysis that shares one
    global weight vector across residues could potentially reduce the explicit
    $d$ factor in the adaptive-regret term.

    \item \textbf{Plug-in versus certified use.}
    In applications with a separate outer calibration controller, the empirical
    CDF may be used directly for efficiency while the lower confidence surface
    is retained as an optional conservative mode.  The same bound quantifies
    the estimation error in either case.
\end{enumerate}

\section{Summary}

For delayed bivariate online observations, the key construction is
\[
    \text{delay staggering}
    \;\Longrightarrow\;
    \text{martingale-difference streams}
    \;\Longrightarrow\;
    \text{uniform finite-grid CDF bound}.
\]
The main edge-level result is
\[
    \sup_{g\in\Q}
    |\widehat p_t(g)-\bar p_t(g)|
    \le
    \sqrt{
        \frac{d}{2N_t}
        \log\!\left(
            \frac{\pi^2 dQ N_t^2}{3\delta}
        \right)
    }
\]
with one global confidence event.  When $N_t\asymp t$, this is
\[
    O\!\left(
        \sqrt{
            \frac{d}{t}
            \log\!\left(\frac{dQt^2}{\delta}\right)
        }
    \right).
\]  Under a stable conditional CDF, this is a
confidence band for the common bivariate CDF.  Under drift, the same statistical
term combines with an explicit tracking-bias term.  A dictionary of moving-window
CDF estimators can additionally be combined with data-dependent fixed-share
weights: the simultaneous confidence event survives post-selection, while the
residue-wise meta-learner competes on every interval with the best memory scale
in Brier loss and in expected grid-averaged squared CDF error.  Under a cumulative
lower-orthant variation budget $V_I$, this yields a best-window contribution of
order $\sqrt{|I|dV_I}+dV_I$ up to constants and the separate meta-regret term.
The construction requires no i.i.d. assumption and treats delay as a filtration problem rather
than as an independence assumption.

\section{Summary of regret guarantees and assumptions}
\label{sec:regret-summary}

This section collects the regret guarantees in one place and separates the
assumptions required for the basic adaptive-regret statement from those needed
only for the variation-budget rate.

Let $I=[r,s]$ be a physical-time interval after burn-in,
$n_I:=|I|$, and let $d_I$ denote the number of residue classes represented in
$I$.  Recall the current-CDF estimation risk
\[
    \mathcal R_k(q)
    :=
    \frac1Q
    \sum_{g\in\Q}
    \bigl(q(g)-p_k(g)\bigr)^2.
\]
Since $\mathcal R_k(p_k)=0$, the quantity
\[
    \sum_{k\in I}
    \E\!\left[
        \mathcal R_k(\widehat p_k^{\mathrm{mix}})
    \right]
\]
is also the cumulative expected excess CDF-estimation risk relative to a
clairvoyant oracle that knows the current conditional CDF $p_k$ at every
origin.

\paragraph{1. Pathwise adaptive regret in Brier loss.}
For the residue-wise fixed-share mixture,
\[
\boxed{
    (1-\eta)
    \sum_{k\in I}
    \ell_k^{\mathrm{mix}}
    \le
    \min_{j\in[J]}
    \sum_{k\in I}
    \ell_{k,j}
    +
    \frac{
        d_I[\log(n_{\max}J)+2]
    }{\eta}.
}
\]
Thus, on every physical-time interval, the adaptive mixture competes with the
best fixed memory scale in the dictionary on that interval.  This is an
adaptive-regret guarantee in the sense of local comparison on arbitrary
contiguous intervals; such interval-wise regret criteria are standard in
adaptive online learning
\cite{herbster1998,hazan2009,gradu2023}.

\paragraph{2. Oracle inequality for current-CDF estimation.}
Using the conditional Brier decomposition gives
\[
\boxed{
    \sum_{k\in I}
    \E\!\left[
        \mathcal R_k(\widehat p_k^{\mathrm{mix}})
    \right]
    \le
    \frac1{1-\eta}
    \left[
        \min_{j\in[J]}
        \sum_{k\in I}
        \E\!\left[
            \mathcal R_k(\widehat p_{k,j})
        \right]
        +
        \frac{
            d_I[\log(n_{\max}J)+2]
        }{\eta}
        +
        \frac{\eta n_I}{4}
    \right].
}
\]
The left-hand side measures distance to the zero-risk truth oracle $p_k$.
The first term on the right is the estimation error of the best available
window in hindsight, while the remaining terms are the additional cost of
adaptively learning the appropriate memory scale.

\paragraph{3. Variation-budget bound.}
Let
\[
    v_u
    :=
    \sup_{g\in\Q}
    |p_{u+1}(g)-p_u(g)|
\]
and
\[
    \mathcal V_{I,j}
    :=
    V_{r-d-W_j+1:s}.
\]
With
\[
    A_j
    :=
    d+\frac{W_j-1}{2},
\]
the adaptive estimator satisfies
\[
\boxed{
\begin{aligned}
    \sum_{k\in I}
    \E\!\left[
        \mathcal R_k(\widehat p_k^{\mathrm{mix}})
    \right]
    \le
    \frac1{1-\eta}
    \Bigg[
        \min_{j\in[J]}
        \left\{
            \frac{n_I d}{2W_j}
            +
            2A_j\E[\mathcal V_{I,j}]
        \right\}
        +
        \frac{
            d_I[\log(n_{\max}J)+2]
        }{\eta}
        +
        \frac{\eta n_I}{4}
    \Bigg].
\end{aligned}
}
\]
The term $n_Id/(2W_j)$ is the statistical fluctuation cost of window $j$,
while $2A_j\E[\mathcal V_{I,j}]$ is its tracking cost under temporal
nonstationarity.

If a common pathwise budget
\[
    \mathcal V_{I,j}\le V_I
    \qquad
    \text{for all }j
\]
is available, the window-dependent part is
\[
    \frac{n_Id}{2W_j}
    +
    2
    \left(
        d+\frac{W_j-1}{2}
    \right)V_I.
\]
Its continuous optimal scale is
\[
    W_I^\star
    =
    \sqrt{\frac{n_Id}{2V_I}},
\]
and a dictionary containing a window within a constant factor of
$W_I^\star$ yields
\[
\boxed{
    O\!\left(
        \sqrt{n_IdV_I}
        +
        dV_I
    \right)
}
\]
for the best-window contribution, in addition to the fixed-share
meta-regret terms.  Variation budgets are a standard way of describing
nonstationarity, although the lower-orthant CDF variation used here differs
from the cost-function variation studied, for example, by
Besbes, Gur, and Zeevi \cite{besbes2015}.

\subsection*{Assumptions used by the regret analysis}

\begin{enumerate}[leftmargin=*]

    \item \textbf{Adapted observations and conditional-CDF target.}
    For every query $g\in\Q$,
    \[
        p_k(g)
        =
        \E[X_k(g)\mid\F_k].
    \]
    The process may otherwise depend arbitrarily on the past.  No i.i.d.,
    mixing, or independence-across-time assumption is imposed.
    Martingale formulations based on conditional distributions under adapted
    dependence are also used in
    \cite{rakhlin2015,yin2022,mineiro2023}.

    \item \textbf{Deterministic bounded reveal delay.}
    Each observation used by the edge is completely revealed after delay $d$.
    The residue construction converts this delayed problem into $d$
    ordinary full-information expert streams.  Online learning with delayed
    feedback and reductions from delayed to non-delayed learning are studied
    more generally by Joulani, Gyorgy, and Szepesvari
    \cite{joulani2013}.
    For a pairwise edge $e=(h_1,h_2)$ in the trajectory application, the
    natural choice is
    \[
        d_e=\max\{d_{h_1},d_{h_2}\}.
    \]

    \item \textbf{Finite query grid.}
    The regret loss is evaluated on a finite set
    $\Q$ of size $Q$:
    \[
        \ell_{k,j}
        =
        \frac1Q
        \sum_{g\in\Q}
        (X_k(g)-\widehat p_{k,j}(g))^2.
    \]
    This is both the computational representation of the bivariate CDF and a
    way of avoiding additional sequential-complexity assumptions required for
    unrestricted continuum-uniform empirical-process results under arbitrary
    dependence; see \cite{rakhlin2015}.  The regret results therefore make no
    claim about a continuum of thresholds outside $\Q$.

    \item \textbf{Finite dictionary of memory scales.}
    The learner competes with a finite collection
    \[
        1\le W_1<\cdots<W_J.
    \]
    Finite expert classes and multiplicative-weight aggregation are standard
    in prediction with expert advice \cite{cesabianchi2006}.  The
    fixed-share construction used to track locally best experts originates
    with Herbster and Warmuth \cite{herbster1998}; our particular
    interval-regret argument follows the fixed-share construction used by
    Gradu, Hazan, and Minasyan \cite{gradu2023}.

    \item \textbf{Predictable mixture weights.}
    The mixture vector $\pi_k$ is $\F_k$-measurable: it may depend on all
    information available before the outcome generated at origin $k$ is
    revealed, but not on that unrevealed outcome.  This is the standard
    non-anticipation condition in online prediction with expert advice
    \cite{cesabianchi2006}.

    \item \textbf{Bounded Brier loss.}
    Every base prediction and mixture prediction takes values in $[0,1]$,
    giving
    \[
        0\le \ell_{k,j}\le1.
    \]
    This boundedness is used in the multiplicative-weight potential argument.
    The Brier score is a classical proper scoring rule for probabilistic
    predictions; see Gneiting and Raftery \cite{gneiting2007}.
    Its conditional decomposition is especially convenient here because the
    excess expected Brier loss relative to the true conditional probability is
    exactly squared estimation error.

    \item \textbf{Fixed-share tuning and finite meta-horizon.}
    The displayed constants use
    \[
        n_{\max}
        =
        \left\lceil
            \frac{T_{\max}}{d}
        \right\rceil,
        \qquad
        \sigma=\frac1{n_{\max}},
        \qquad
        0<\eta\le\frac12.
    \]
    These are algorithmic tuning conditions rather than distributional
    assumptions.  Fixed-share and related expert-tracking methods use
    analogous sharing parameters
    \cite{herbster1998,gradu2023}.

    \item \textbf{Variation is required only for the explicit
    truth-oracle rate.}
    The fixed-share adaptive-regret and local oracle inequalities do
    \emph{not} assume a drift bound, $\rho$, or a variation budget.
    The quantities $\mathcal V_{I,j}$ and $V_I$ enter only when converting
    the best-window risk into an explicit rate relative to the current true
    CDF.  Variation-budget models for nonstationary environments are standard
    in the nonstationary optimization literature
    \cite{besbes2015}; the variation metric used in this note is specific to
    lower-orthant conditional CDFs.

    \item \textbf{Near-oracle dictionary coverage is required only for the
    simplified rate.}
    The theorem itself holds for any finite dictionary
    $\{W_1,\ldots,W_J\}$.  The rate
    \[
        O\!\left(
            \sqrt{n_IdV_I}+dV_I
        \right)
    \]
    additionally assumes that the dictionary contains a memory length within
    a constant factor of $W_I^\star$.  A geometric dictionary
    $W_j=2^{j-1}$ provides this coverage away from the admissible endpoints.
    This is a design choice, not an assumption on the data-generating process.

    \item \textbf{Burn-in.}
    The regret statements are written after the largest required base window
    is available.  This is a bookkeeping condition ensuring that every expert
    in the dictionary is defined, rather than a stochastic assumption.

\end{enumerate}

\paragraph{What is not assumed.}
In particular, the regret analysis does not require independence, stationarity,
mixing, a known drift rate $\rho$, a known variation budget $V_I$, or a
confidence level $\delta$.  The parameter $\delta$ belongs to the separate
high-probability concentration analysis; it does not enter the fixed-share
regret or variation-budget oracle inequalities.


\begin{thebibliography}{99}

\bibitem{rakhlin2015}
A. Rakhlin, K. Sridharan, and A. Tewari.
\newblock Sequential complexities and uniform martingale laws of large numbers.
\newblock \emph{Probability Theory and Related Fields},
161:111--153, 2015.
\newblock
\href{https://www.mit.edu/~rakhlin/papers/emp_proc_dep_revised.pdf}
{Preprint}.

\bibitem{yin2022}
S. Yin, S. Agrawal, and A. Zeevi.
\newblock Online allocation and learning in the presence of strategic agents.
\newblock \emph{Advances in Neural Information Processing Systems 35}, 2022.
\newblock
\href{https://papers.neurips.cc/paper_files/paper/2022/file/29455889cb812d661573f6fbc800bd9e-Paper-Conference.pdf}
{Paper}.

\bibitem{mineiro2023}
P. Mineiro and S. Howard.
\newblock Time-uniform confidence bands for the CDF under nonstationarity.
\newblock \emph{Advances in Neural Information Processing Systems 36}, 2023.
\newblock
\href{https://proceedings.neurips.cc/paper_files/paper/2023/hash/148bbc25b934211d80435b5cad5a7198-Abstract-Conference.html}
{Paper}.

\bibitem{cesabianchi2006}
N. Cesa-Bianchi and G. Lugosi.
\newblock \emph{Prediction, Learning, and Games}.
\newblock Cambridge University Press, 2006.
\newblock
\href{https://doi.org/10.1017/CBO9780511546921}
{Book}.

\bibitem{herbster1998}
M. Herbster and M. K. Warmuth.
\newblock Tracking the best expert.
\newblock \emph{Machine Learning}, 32(2):151--178, 1998.
\newblock
\href{https://doi.org/10.1023/A:1007424614876}
{Paper}.

\bibitem{hazan2009}
E. Hazan and C. Seshadhri.
\newblock Efficient learning algorithms for changing environments.
\newblock In \emph{Proceedings of the 26th International Conference on
Machine Learning}, pages 393--400, 2009.
\newblock
\href{https://doi.org/10.1145/1553374.1553425}
{Paper}.

\bibitem{gradu2023}
P. Gradu, E. Hazan, and E. Minasyan.
\newblock Adaptive regret for control of time-varying dynamics.
\newblock \emph{Proceedings of Machine Learning Research},
211:560--572, 2023.
\newblock The fixed-share meta-algorithm and interval-regret proof used here
appear as Algorithm~2 and Lemma~A.2 in the extended version,
\href{https://arxiv.org/pdf/2007.04393}
{arXiv:2007.04393v3}.

\bibitem{joulani2013}
P. Joulani, A. Gyorgy, and C. Szepesvari.
\newblock Online learning under delayed feedback.
\newblock In \emph{Proceedings of the 30th International Conference on
Machine Learning},
\emph{Proceedings of Machine Learning Research} 28(3):1453--1461, 2013.
\newblock
\href{https://proceedings.mlr.press/v28/joulani13.html}
{Paper}.

\bibitem{gneiting2007}
T. Gneiting and A. E. Raftery.
\newblock Strictly proper scoring rules, prediction, and estimation.
\newblock \emph{Journal of the American Statistical Association},
102(477):359--378, 2007.
\newblock
\href{https://doi.org/10.1198/016214506000001437}
{Paper}.

\bibitem{besbes2015}
O. Besbes, Y. Gur, and A. Zeevi.
\newblock Non-stationary stochastic optimization.
\newblock \emph{Operations Research}, 63(5):1227--1244, 2015.
\newblock
\href{https://doi.org/10.1287/opre.2015.1408}
{Paper}.

\bibitem{massart1990}
P. Massart.
\newblock The tight constant in the Dvoretzky--Kiefer--Wolfowitz inequality.
\newblock \emph{The Annals of Probability}, 18(3):1269--1283, 1990.

\end{thebibliography}

\end{document}
